\documentclass{article}
\usepackage[utf8]{inputenc}
\usepackage{amsmath}
\usepackage{amssymb}
\usepackage{geometry}
\usepackage{hyperref}
\geometry{a4paper, margin=1in}

\title{Supplementary Materials: Mathematical Derivations and Lean 4 Formalizations}
\author{}
\date{}

\begin{document}
\maketitle

\section*{Overview}
This supplementary document provides the formal mathematical framework underlying the core claims presented in the main manuscript. The theorems below have been formalized and verified using the Lean 4 interactive theorem prover. Formalization ensures that the transitions from foundational physical primitives (symmetries, topology, and thermodynamics) to the complex structural properties of consciousness (boundaries, unity, and reflexivity) are stated precisely rather than metaphorically. It does not by itself make them physically true. The development now declares \emph{no} axioms --- \texttt{\#print axioms} reports only Lean's three foundational ones --- but the physical postulates have not vanished; they have moved into class fields that each instance must discharge, and into the hypotheses of individual theorems. Five axioms in earlier drafts were removed after three were found to be logically inconsistent; see the soundness subsection of the main text. Where the Lean statement is weaker than the informal theorem above it, we say so explicitly in the implementation note. The complete Lean 4 source code is publicly available at \url{https://github.com/bobaseb/physical_primitives_of_mind}, and web-accessible documentation is available at \url{https://bobaseb.github.io/physical_primitives_of_mind/}.

\section{Symmetry Breaking and the Emergence of Boundaries}
The first step in our deductive chain establishes that a perfectly symmetric universe cannot support agency. Any distinguishable physical system requires a topological defect.

Let $\mathcal{M}$ be a pseudo-Riemannian manifold representing spacetime, and let $\phi: \mathcal{M} \to \mathcal{V}$ be a physical field mapping to a vacuum manifold $\mathcal{V}$ (the space of energy minima). Let $G$ be a continuous symmetry group acting on $\mathcal{V}$. 

\textbf{Theorem 1 (Topological Defects from Broken Symmetry):} If the continuous symmetry $G$ is spontaneously broken such that $\mathcal{V} \cong G/H$ (where $H$ is the unbroken subgroup), and the fundamental group $\pi_1(\mathcal{V}) \neq 0$, then field configurations $\phi$ defined on a spatial boundary $S^1 \subset \mathcal{M}$ cannot necessarily be smoothly deformed to a constant vacuum state.

\textit{Lean 4 Implementation (\texttt{Phase1\_Primitives.lean}):} We formalize this by defining a \texttt{VacuumManifold} and proving \texttt{boundary\_defect\_forces\_interior\_vacuum\_break}: if a boundary field configuration $\phi_{\partial}$ is not null-homotopic, then it admits no continuous extension over a contractible interior. This geometric obstruction forms the fundamental physical boundary separating an ``inside'' from an ``outside.'' Note the direction: Lean establishes that a defect \emph{obstructs extension}, not that a defect must exist. The companion lemma \texttt{contractible\_interior\_forces\_trivial\_boundary} (renamed from \texttt{defect\_inevitability}, whose old name asserted the opposite of its statement) proves the contrapositive. Separately, \texttt{spontaneous\_symmetry\_breaking} shows a global energy minimiser lies in the vacuum manifold almost everywhere; this is now axiom-free, the pointwise step being \texttt{pointwise\_vacuum\_of\_global\_min}.

\section{The Thermodynamic Constraint: Landauer's Principle}
Once a boundary is established, it is subjected to external perturbations. Because the internal phase space of any localized physical system is finite, the system must overwrite internal states to adapt.

Let the global phase space $\Omega = \Omega_{sys} \times \Omega_{bath}$ undergo globally reversible Hamiltonian dynamics, modeled as an injective mapping $f: \Omega \to \Omega$.

\textbf{Theorem 2 (Thermodynamic Erasure):} If the local transition map on the system $\pi_{sys} \circ f$ is not injective (i.e., multiple initial states map to the same final state, contracting the local phase space), then by the conservation of global phase space volume (Liouville's theorem), the bath phase space must expand, resulting in an irreversible transfer of heat to the environment: $\Delta Q \ge k_B T \ln 2$.

\textit{Lean 4 Implementation (\texttt{Phase3\_CombinatorialThermodynamics.lean}):} We explicitly decouple logical erasure from physical erasure. Lean verifies that any contraction in the \textit{physical} state space ($\neg \text{Function.Injective}(t_{sys})$) under globally reversible dynamics mathematically necessitates heat dissipation. The combinatorial core, \texttt{landauer\_from\_reversibility} and \texttt{second\_law}, is fully derived. Converting it into a statement about heat requires the postulate $\Delta Q = T \Delta S_{bath}$, which is carried as the class field \texttt{StatisticalMechanics.heat\_eq}; \texttt{Examples.lean} exhibits a one-bit erasure system with an explicit bath that discharges it, so the postulate is demonstrably satisfiable. The boundary is thus forced to operate as a dissipative structure.

\section{Structural Resonance and Least Action}
To avoid catastrophic dissipation, the dissipative structure evolves to minimize entropy production. Let $P(ext)$ be the probability measure of external environmental perturbations, and $Q(int)$ be the internal transition probabilities of the system.

The total entropy production $\sigma$ is bounded from below by the Kullback-Leibler divergence:
\begin{equation}
\sigma \ge \frac{1}{\Delta t} D_{KL}(P(ext) || Q(int))
\end{equation}

\textbf{Theorem 3 (Structural Resonance):} By the Principle of Least Action, the system undergoes gradient descent on $\sigma$, forcing $\lim_{t \to \infty} D_{KL}(P || Q) = 0$.

This conceptual mapping suggests that ``prediction error minimization'' is fundamentally a physical thermodynamic attractor. The information-theoretic KL bound is formalized in \texttt{Phase3\_KLBound.lean} as the class field \texttt{StructuralResonance.kl\_bound} --- an irreducible physical postulate attached to a specific system rather than asserted globally, the earlier \texttt{kl\_bound\_axiom} formulation having been inconsistent. Gibbs' inequality (KL non-negativity) is proved as a genuine theorem from $\log x \le x - 1$; \texttt{structural\_resonance\_bound} restates the axiom in the form $D_{KL} \le \Delta t \cdot \sigma$, and \texttt{discrete\_entropy\_rate\_nonneg} combines the two to derive $\sigma \ge 0$. The gradient-flow half of Theorem 3 is now formalized against the entropy production functional itself: \texttt{Phase8\_ContinuousField.lean} \S7 derives the Fr\'echet derivative of $\sigma$ with respect to the coupling kernel, identifies the differentiated functional with \texttt{entropy\_production\_rate} on finite substrates carrying counting measure, and concludes in \texttt{structural\_resonance\_decreases\_entropy\_production} that $\sigma$ is non-increasing along the resulting flow; \texttt{Examples.lean} \S5 exhibits a two-site substrate on which $\sigma$ decays as $e^{-2t}$. What remains unformalized is the limit itself, $\lim_{t\to\infty} D_{KL} = 0$: monotone descent plus the bound $D_{KL} \le \Delta t\,\sigma$ does not by itself force the limit to be zero, which would need a coercivity or {\L}ojasiewicz-type estimate we do not have. The link is also restricted to finite substrates and to a fixed phase field.

\section{Macroscopic Scaling and Kuramoto Phase Synchronization}
As the system scales, it forms a continuous macroscopic electrodynamic field. We model the local electromagnetic phases $\theta_i$ as coupled non-linear oscillators via the Kuramoto model. 

For a system with coupling strength $K$ and thermodynamic noise $D$, the dynamic evolution is:
\begin{equation}
\dot{\theta_i} = \omega_i + \frac{K}{N} \sum_{j=1}^N \sin(\theta_j - \theta_i) + \eta_i(t)
\end{equation}

\textbf{Theorem 4 (Critical Phase Transition):} When the integrated spatial coupling $K$ exceeds the thermodynamic noise threshold $K_c = 2D$, the oscillators undergo a sudden phase transition into a synchronized, phase-locked state ($r > 0$).

\textit{Lean 4 Implementation (\texttt{Phase4\_KuramotoDynamics.lean}):} We formalize the Lyapunov potential $V(\theta) = -\tfrac{1}{2}\sum_{i,j} A_{ij}\cos(\theta_j - \theta_i)$ of the coupled field and prove two directions: \texttt{phase\_locked\_minimizes\_potential} (the phase-locked configuration attains the global minimum whenever all couplings are positive) and \texttt{potential\_min\_implies\_phase\_locked} (any minimizer is phase-locked). Note the scope: this is a statement about the \emph{static} potential landscape and holds for all positive coupling matrices. The critical threshold $K_c = 2D$ does not appear in this file and is not formalized in Lean --- establishing it requires the noise-driven stochastic dynamics and the spectral theory of the associated Fokker-Planck operator, neither of which is currently available in Mathlib. The threshold is instead validated numerically (\texttt{simulations/kuramoto.py}); see Table 1 of the main text, where it is listed with status ``Numerical.''

\textit{Lean 4 Implementation (\texttt{Phase8\_ContinuousField.lean}):} What the continuous-field file does carry is the threshold value itself, \texttt{critical\_coupling D = 2 * D}, and a predicate \texttt{exhibits\_phase\_transition} for a substrate lying above it. Two cautions. First, the predicate is a definition with no theorem connecting it to \texttt{is\_continuous\_kuramoto\_trajectory}, \texttt{entropy\_production\_rate}, or any order parameter: $2D$ is stipulated in the source, and its justification is the mean-field self-consistency equation $r = I_1(Kr/D)/I_0(Kr/D)$, whose bifurcation at $K = 2D$ we do not formalize. Second, the predicate now requires the substrate's \texttt{volume} to be a probability measure, and this hypothesis is load-bearing rather than decorative. It compares \texttt{mean\_field\_coupling}, the kernel averaged over both arguments, against $2D$; on a substrate of total mass $m$ that average carries a factor $m^2$, so without normalization the comparison tracks the substrate's size rather than its coupling and can be satisfied by any kernel at all. \texttt{Examples.lean}~\S9 exhibits a system a factor of two below threshold whose unnormalized double integral nonetheless clears $2D$. Under the hypothesis, \texttt{mean\_field\_coupling\_const} gives that a constant kernel has strength exactly $K$ and \texttt{exhibits\_phase\_transition\_const\_iff} reduces the predicate to $K > 2D$, witnessed in \S9 on both sides of the threshold.

\textit{Lean 4 Implementation (\texttt{Phase4\_RotatingFrame.lean}):} The potential above omits the natural-frequency term $-\sum_i \omega_i\theta_i$ that appears in the Lyapunov descent result \texttt{dV\_dt\_le\_zero} of \texttt{Phase3\_CombinatorialThermodynamics.lean}. The omission is not cosmetic. With any $\omega_i \ne 0$ the full potential is unbounded below (\texttt{kuramoto\_potential\_unbounded\_below}), so it has no minimum for a phase-locked state to attain, and the two results concerned different functionals. This file chains them by the standard rotating-frame reduction $\theta_i \mapsto \theta_i - \Omega t$. For a system of identical natural frequencies $\omega \equiv \Omega$: \texttt{is\_kuramoto\_trajectory\_rotate} sends a trajectory to a trajectory of the zero-frequency system \texttt{sys.reduced} --- the same system Theorem 5's \texttt{ThermodynamicCover} works with; \texttt{kuramoto\_potential\_reduced} shows the two potentials coincide on it; \texttt{dynamic\_potential\_descent} transports the descent identity $\dot V = -\sum_i \dot\theta_i^2$ onto the potential above; and \texttt{rotating\_frame\_chain} runs this through the minimum characterization to phase-locking and $r^2 = 1$. Phase differences, phase-locking and the order-parameter magnitude are proved frame-invariant (\texttt{rotate\_sub}, \texttt{is\_phase\_locked\_rotate}, \texttt{order\_parameter\_r\_sq\_shift}), so no observable is lost. \texttt{Examples.lean}~\S7 runs the chain end to end on a two-oscillator system, discharging every hypothesis, and exhibits a system whose full potential has no minimum. Scope: the reduction is exact only for identical frequencies. A spread of frequencies leaves residual detunings $\omega_i - \Omega$ in the reduced system, to which the unboundedness result still applies, and phase-locking then depends on $K_c$, which is not formalized.

\textit{From the continuous field to the discrete model (\texttt{Phase2\_SimplicialBridge.lean}, \texttt{Phase2\_MeshConvergence.lean}):} The coupling matrix $A_{ij}$ above is not postulated. A triangulation of the substrate assigns to each edge $(u,v)$ a region of the manifold, and the coupling weight is the Lebesgue--Bochner integral of the scalar magnitude of the stress-energy tensor over that region (\texttt{edge\_weight}); symmetry and non-negativity of $A$ are then proved rather than assumed. The bridge back is mesh refinement: as the regions shrink, the discrete energy $\tfrac12\sum_{u,v} \mu(R_{uv})\,f(x_u)$ should converge to the continuous energy $\int_S f$. This was previously stated as a conjecture and was in fact false as stated --- nothing forced the regions to cover the manifold, and empty regions have diameter zero, so the everywhere-empty triangulation refuted it. It is now a theorem (\texttt{mesh\_refinement\_convergence}), quantified over a sequence of triangulations of a region $S$ whose regions are measurable, pairwise disjoint across distinct edges, covering, and anchored to the embedded vertices, with mesh size tending to zero. The fixed-mesh error is bounded by $\varepsilon\,\mu(S)$, with $\varepsilon$ a modulus of continuity of the integrand, and the symmetric double sum is proved equal to the midpoint rule over unordered edges. Two limitations are worth stating. The theorem gives convergence, not the $O(1/N^2)$ rate reported by \texttt{simulations/mesh\_refinement.py}. And \texttt{TriangulatedManifold} itself remains an unconstrained class: the geometric conditions live in the predicate \texttt{IsRegularTriangulation}, so results that need them say so in their hypotheses. Because those conditions are in tension --- cells disjoint yet covering, every vertex anchored to each of its edges --- \texttt{Examples.lean}~\S6 exhibits an instance: the uniform partition of $[0,1)$ into $N$ half-open cells, on which the convergence theorem is instantiated and the approximations are checked to genuinely change with $N$.

\section{Global Sections and the Unity of Consciousness}
Phase synchronization alone is mathematically insufficient for conscious unity; the synchronized local fields must be topologically glued together.

\textbf{Theorem 5 (Sheaf-Theoretic Unity):} Given a presheaf of probability measures $\mathcal{F}$ over a continuous biological substrate $X$, if local invariant measures $P_U \in \mathcal{F}(U)$ agree on all spatial intersections $U \cap V$ (guaranteed by the Kuramoto phase-locking), they uniquely glue into a singular Global Section $P_X \in \mathcal{F}(X)$.

\textit{Lean 4 Implementation (\texttt{Phase5\_GlobalSection.lean}):} We apply sheaf theory to prove that phase-locked local fields fuse into a single global section, via \texttt{global\_section\_from\_thermodynamics}. One caveat and one discharged obligation. The caveat: phase-locking is not derived from dynamics here, it is the class field \texttt{ThermodynamicCover.thermodynamic\_equilibrium}, so the theorem reads ``given a cover already at the potential minimum, the local sections glue uniquely.'' The discharged obligation: an instance of \texttt{ThermodynamicCover} now exists (\texttt{Examples.lean}, \S4), so the theorem is no longer conditional on a structure of unknown realizability. The witness is a three-site substrate with the discrete topology and its Borel $\sigma$-algebra, covered by two overlapping open patches meeting in one shared site; local sections are germ families of the phase-indexed measure $\mu_\theta = (1+\cos\theta)^{+}\,\delta_{\mathrm{mid}}$ in the sheafification of the probability presheaf. Two details matter for non-vacuity: $\mu_\theta$ is provably non-constant in $\theta$ (\texttt{phaseMeasure\_not\_const}), so \texttt{phase\_invariant\_periodic} is a real $2\pi$-periodicity obligation rather than a triviality; and \texttt{thermodynamic\_equilibrium} is discharged from \texttt{phase\_locked\_minimizes\_potential} rather than assumed. We therefore claim a formalized, non-vacuous gluing mechanism --- not a solved binding problem, since the witness is a toy substrate whose phases are locked by construction.

\section{Reflexive Topology and the Self}
To minimize internal thermodynamic friction, the unified Global Section must model its own internal state changes.

\textbf{Theorem 6 (Auto-Resonance):} The system minimizes total entropy production only if the internal state mapping $Q(int)$ contains a low-dimensional reflexive projection of its own boundary $\partial \mathcal{M}$.

\textit{Lean 4 Implementation (\texttt{Phase6\_ReflexiveTopology.lean}):} We formalize the reflexive projection as a contraction on a complete metric space and apply the Banach fixed-point theorem to obtain a unique fixed point --- the ``Self'' as a topological feature rather than a metaphysical addition. The entropy production functional does not appear in this file: the identification of the Banach fixed point with the minimizer of $\sigma$ is an interpretive step made in the main text, and the contraction constant ($1/2$) is assumed rather than derived from field dynamics.

\section{Topological Rigidity of von Neumann Architectures}
Finally, we compare the continuous coupling matrix $A(x,y)$ of biological fields with the discrete, rigid routing matrix $G_{ij}$ of standard artificial hardware (e.g., GPUs).

\textbf{Theorem 7 (Hardware Thermodynamics):} Because a continuous parameter space contains global minima inaccessible to a discrete lattice under identical thermodynamic constraints, the dynamical potential of a biological neural field is strictly lower than that of rigid silicon hardware: $V(continuous) < V(discrete)$.

\textit{Lean 4 Implementation (\texttt{Phase7\_Hardware\-Comparison.lean}):} The lemma \texttt{exists\_better\_coupling\_allocation}, wrapped as the theorem \texttt{continuous\_beats\_rigid\_topology}, proves the reallocation step: given a coupling matrix satisfying \texttt{is\_strictly\_suboptimal} (positive weight on a less phase-correlated pair while a more correlated pair exists), there is another valid coupling with identical total resources and strictly greater total phase correlation.

\textit{Lean 4 Implementation (\texttt{Phase7\_Rigidity.lean}):} Two defects of the above are addressed here. \emph{(i) Rigidity as a support constraint.} \texttt{RealizableIn S R A} says $A$ is a valid coupling meeting resource budget $R$ and vanishing on every pair the architecture has no wire for --- weights tunable, wiring fixed. \texttt{totalCorrelation\_le\_of\_realizable} proves no realizable coupling exceeds $R\cdot M$ when $M$ bounds phase correlation over the wired pairs, and \texttt{exists\_realizable\_pair} attains it. \texttt{rigid\_is\_strictly\_suboptimal} then \emph{derives} \texttt{is\_strictly\_suboptimal} for every realizable coupling, whenever some pair $(k,l)$ the wiring misses has correlation exceeding $M$: the hypothesis Theorem 7 previously assumed is now a consequence of the wiring pattern. \texttt{rigid\_gap} gives the quantitative form, a gain of at least $R\,(\cos(\theta_l-\theta_k) - M) > 0$ at equal budget. \texttt{Examples.lean}~\S8 witnesses this on three sites wired only to the maximally anti-correlated pair, where the gap is exactly $2$, and runs \texttt{exists\_better\_coupling\_allocation} on it with no assumed hypothesis. \emph{(ii) Continuity as a measure-theoretic distinction.} \texttt{fieldCorrelation} integrates a kernel against $\mu\otimes\mu$ as in Derivation 4. A finitely-sited architecture is a \texttt{Finset X}; \texttt{fieldCorrelation\_sited\_eq\_zero} proves that on an atomless substrate such a kernel contributes exactly zero, for every phase field and every choice of weights, while \texttt{fieldCorrelation\_patchKernel} gives $c\,\mu(U)^2$ for a patch. \texttt{sited\_architecture\_below\_field\_optimum} combines them; \texttt{Examples.lean}~\S8 instantiates it on $\mathbb{R}$ with a thousand units of weight at the origin against the unit interval.

\textbf{Scope caveat.} Result (i) is resource-matched and no longer question-begging, but it turns on fixed wiring support, not on continuity: an architecture wired to the best pair is not beaten. Result (ii) is a genuine type-level separation but is \emph{not} resource-matched, since a measure-zero substrate carries zero resource in the $\mu\otimes\mu$ sense as well; what it establishes is that the continuum coupling energy cannot register a finitely-sited architecture at all, which is a claim about the model's observable rather than about computational power. Neither result formalizes manifold structure or lattice geometry. Accordingly, the inference to ``von Neumann architectures cannot experience unified consciousness'' remains an informal argument in the main text, now with two verified steps rather than one.

\end{document}
